\section*{Bab \ref{ch:b2:population-genetics} --- Genetika Populasi}

\begin{solution}{exo:b2:population-genetics:1}
$p(M) = (1280 + 320)/2000 = 0.8$, $q(N) = 0.2$. Harapan $MM$:
$0.64\times 1000 = 640$; $MN$: $2\times 0.8\times 0.2\times 1000 =
320$; $NN$: 40 --- tepat sama dengan jumlah yang teramati. Populasinya
berada dalam kesetimbangan pada lokus ini.
\end{solution}

\begin{solution}{exo:b2:population-genetics:2}
$q = \sqrt{1/20000} = 0.0071$; pembawanya $2pq \approx 0.014$, yaitu
satu orang dari 70. Pembawanya melebihi orang yang terkena sebanyak
$2p/q \approx
280$ berbanding satu.
\end{solution}

\begin{solution}{exo:b2:population-genetics:3}
\omterm{def:b2:population-genetics:fitness}{Kebugaran}: jumlah keturunan harapan sebuah \omterm{def:b2:meiosis-heredity:vocabulary}{genotipe} terhadap yang
terbaik. \omterm{def:b2:population-genetics:fitness}{Koefisien seleksi} $s$: kekurangan \omterm{def:b2:population-genetics:fitness}{kebugarannya}, $w = 1 -
s$.
Heritabilitas $h^{2}$: bagian ragam \omterm{def:b2:meiosis-heredity:vocabulary}{fenotipe} yang bersifat genetik
aditif, yang setara dengan kemiringan keturunan terhadap induknya.
Gayanya: \omterm{def:b2:genome-diversification:mutation}{mutasi}, seleksi, hanyutan, migrasi, dan perkawinan tak acak
(yang terakhir mengubah frekuensi \omterm{def:b2:meiosis-heredity:vocabulary}{genotipe}, bukan \omterm{def:b2:population-genetics:frequencies}{frekuensi alel}).
\end{solution}

\begin{solution}{exo:b2:population-genetics:4}
Berarah: melanisme pada ngengat berbedak, dengan agen berupa burung
yang berburu dengan penglihatan pada batang yang menggelap oleh jelaga.
Pemantap: bobot lahir manusia, dengan agen berupa kematian bayi yang
sangat kecil dan yang sangat besar. Penyeimbang: sel sabit, dengan agen
berupa malaria (terhadap $AA$) dan anemia (terhadap $SS$) bersama-sama.
\end{solution}

\begin{solution}{exo:b2:population-genetics:5}
$q_t = q_0/(1 + tq_0)$: menyetengahkannya memerlukan $1/q_0 = 50$
generasi; mencapai $0.001$ memerlukan $1/0.001 - 1/0.02 = 950$
generasi. Hampir seluruh salinan \omterm{def:b2:meiosis-heredity:vocabulary}{alelnya} duduk di dalam \omterm{def:b2:meiosis-heredity:vocabulary}{heterozigot},
yang tak terlihat oleh seleksi; mencegah yang terkena untuk berbiak
(yang sudah dikerjakan alam, karena mereka mati) tidak mengubah apa pun
yang terukur dalam satu hayat manusia.
\end{solution}

\begin{solution}{exo:b2:population-genetics:6}
$p = 0.3$: $\bar w = 0.09 + 0.42 + 0.8\times 0.49 = 0.902$; $p' =
(0.09 + 0.21)/0.902 = 0.333$; $\Delta p = +0.033$. $p = 0.9$: $\bar w
= 0.998$, $p' = 0.9018$, $\Delta p = +0.0018$. Pada $p = 0.9$ hanya
$q^{2} = 0.01$ populasinya yang bergenotipe $aa$: \omterm{def:b2:meiosis-heredity:vocabulary}{alel} yang dikenai
seleksi tersembunyi di dalam \omterm{def:b2:meiosis-heredity:vocabulary}{heterozigot}, dan $\Delta p = spq^{2}/\bar w$ membawa faktor $q^{2}$.
\end{solution}

\begin{solution}{exo:b2:population-genetics:7}
Dengan $s = 0.12$ terhadap $AA$ dan $t = 0.8$ terhadap $SS$, $\hat q =
s/(s + t) = 0.12/0.92 = 0.13$. Kelahiran yang terkena $\hat q^{2} = 0.017$, yaitu kira-kira satu anak dari 60 --- harga perlindungan yang
dinikmati \omterm{def:b2:meiosis-heredity:vocabulary}{heterozigotnya}.
\end{solution}

\begin{solution}{exo:b2:population-genetics:8}
\omterm{def:b2:meiosis-heredity:vocabulary}{Alel} \omterm{def:b2:meiosis-heredity:vocabulary}{resesif} yang mematikan: $\hat q = \sqrt{\mu} = \sqrt{2\times 10^{-5}} =
0.0045$; kelahiran yang terkena $\hat q^{2} = \mu = 2\times 10^{-5}$,
yaitu satu dari \num{50000}. \omterm{def:b2:meiosis-heredity:vocabulary}{Alel} \omterm{def:b2:meiosis-heredity:vocabulary}{dominan} yang mematikan: tiap
salinannya disingkirkan pada generasi kemunculannya, sehingga kelahiran
yang terkena hanyalah \omterm{def:b2:genome-diversification:mutation}{mutasi} yang baru, $2\mu = 4\times 10^{-5}$ (dua
gamet tiap kelahiran), yaitu satu dari \num{25000} --- yang \omterm{def:b2:meiosis-heredity:vocabulary}{resesif}
menyembunyikan \omterm{def:b2:meiosis-heredity:vocabulary}{alelnya} dua ratus berbanding satu, yang \omterm{def:b2:meiosis-heredity:vocabulary}{dominan} tidak
menyembunyikan apa pun.
\end{solution}

\begin{solution}{exo:b2:population-genetics:9}
$H_t = 0.5\,(1 - 1/100)^{t}$: setelah 10 generasi $0.45$; setelah 50,
$0.30$; setelah 200, $0.067$. Bagi $(1 - 1/2N)^{100} = 0.9$:
$100/(2N) \approx -\ln 0.9 = 0.105$, sehingga $N \approx 475$, yaitu
kira-kira 500 individu yang berbiak.
\end{solution}

\begin{solution}{exo:b2:population-genetics:10}
Dua puluh salinan gen: $P(\text{tak ada}) = 0.9^{20} = 0.12$. Sebuah
salinan tunggal di antara 20 berfrekuensi $0.05$, dan sebuah \omterm{def:b2:meiosis-heredity:vocabulary}{alel}
netral terpancang dengan peluang yang sama dengan frekuensinya:
$0.05$. Pada sebuah daratan berisi, katakanlah, \num{100000} burung
sebuah salinan tunggal memiliki peluang $5\times 10^{-6}$ untuk
terpancang. Pulaunya adalah tempat \omterm{def:b2:meiosis-heredity:vocabulary}{alel} yang jarang hilang dan tempat
\omterm{def:b2:meiosis-heredity:vocabulary}{alel} yang bertahan dapat mengambil alih: hanyutannya memiskinkan
sekaligus membedakan.
\end{solution}

\begin{solution}{exo:b2:population-genetics:11}
$R = h^{2}S = 0.3\times 1.4\times \qty{1000}{L} = \qty{420}{L}$ tiap
generasi. Dengan pejantan dari \qty{1}{\%} teratas, diferensial
reratanya adalah $(1.4 + 2.7)/2 = 2.05$ simpangan baku dan $R =
0.3\times 2050 = \qty{615}{L}$. Tanggapannya melambat karena seleksinya
menghabiskan ragam aditifnya --- \omterm{def:b2:meiosis-heredity:vocabulary}{alel} yang menguntungkan pergi menuju
fiksasi dan $h^{2}$ turun --- dan karena perubahan berkorelasi pada
sifat lainnya (kesuburan, kesehatan) membebankan ongkos yang tidak
dilihat seleksi atas hasilnya saja.
\end{solution}

\begin{solution}{exo:b2:population-genetics:12}
Seleksi melihat \omterm{def:b2:meiosis-heredity:vocabulary}{fenotipe}. Sebuah \omterm{def:b2:meiosis-heredity:vocabulary}{alel} \omterm{def:b2:meiosis-heredity:vocabulary}{resesif} di dalam sebuah
\omterm{def:b2:meiosis-heredity:vocabulary}{heterozigot} tidak membuat \omterm{def:b2:meiosis-heredity:vocabulary}{fenotipe}, sehingga seleksi tidak dapat
menyingkirkannya, dan frekuensinya turun sebagai $1/t$, makin lama
makin lambat. Sebuah \omterm{def:b2:genome-diversification:mutation}{mutasi} netral tidak membuat perbedaan \omterm{def:b2:population-genetics:fitness}{kebugaran};
nasibnya hanya milik hanyutan, dan ia terpancang dengan peluang $1/2N$.
Sebuah sifat yang heritabilitasnya nol boleh jadi terseleksi dengan
kuat --- yang bertahan hidup boleh jadi sangat berbeda dari
populasinya --- namun generasi berikutnya tidak berubah, karena
perbedaannya tidak terwariskan. Seleksi hanya dapat mengenai perbedaan
yang terwariskan yang membuat perbedaan bagi \omterm{def:b2:population-genetics:fitness}{kebugarannya}; segala yang
lain berevolusi secara kebetulan atau tidak sama sekali.
\end{solution}

\begin{solution}{pb:b2:population-genetics:1}
\textbf{1.} Bagian melaniknya $1 - q^{2} = 1 - 0.999^{2} = 0.002$,
yaitu satu ngengat dari 500; bagian \omterm{def:b2:meiosis-heredity:vocabulary}{alel} $M$ di dalam \omterm{def:b2:meiosis-heredity:vocabulary}{heterozigot}
adalah $2pq/(2pq + 2p^{2}) = q = 0.999$: pada dasarnya semuanya.
\textbf{2.} Dengan $\bar w = 1 - sq^{2}$, salinan \omterm{def:b2:meiosis-heredity:vocabulary}{alel} yang pucat
setelah seleksinya adalah $pq\cdot 1 + q^{2}(1 - s)$ dibagi $\bar w$, sehingga $q'
= (pq + q^{2} - sq^{2})/(1 - sq^{2}) = q(1 - sq)/(1 - sq^{2})$.
\textbf{3.} $q_1 = 0.999\times(1 - 0.2997)/(1 - 0.2994) = 0.99857$:
sebuah perubahan sebesar $-0.0014$.
\textbf{4.} $q = \sqrt{0.1} = 0.32$.
\textbf{5.} Rekursinya, jika diulang, memerlukan 28 generasi bagi $s =
0.3$; hampirannya memberikan orde yang sama (perubahannya adalah
$-0.0003$ pada mulanya, lalu memercepat ketika $p$ bertambah). Lima
puluh generasi yang teramati bersesuaian dengan $s$ yang lebih kecil,
kira-kira $0.2$ (rekursinya memberikan 45 generasi bagi $s = 0.2$):
seleksi sebesar \qtyrange{20}{30}{\%} tiap generasi, yang luar biasa
menurut ukuran kebanyakan \omterm{def:b2:meiosis-heredity:vocabulary}{alel}.
\textbf{6.} Dengan $w_{MM} = w_{Mm} = 1 - s$, $w_{mm} = 1$: $\bar w = 1
- s(1 - q^{2})$ dan $p' = p(1 - s)/\bar w$, sehingga $\Delta p = p[(1 - s) -
\bar w]/\bar w = -spq^{2}/\bar w$. Bentuk yang pucat hanya disukai
sebagai \omterm{def:b2:meiosis-heredity:vocabulary}{homozigot} $mm$, yaitu sebagian sebesar $q^{2} = 0.1$ pada
awalnya: kembalinya bermula dengan lambat. Ketika $q \to 1$, $\Delta p \approx -sp$ dan $M$ menurun secara eksponensial, dengan faktor $1 - s$
tiap generasi: sebuah \omterm{def:b2:meiosis-heredity:vocabulary}{alel} \omterm{def:b2:meiosis-heredity:vocabulary}{dominan} tidak memiliki tempat bersembunyi.
Kenaikannya adalah bayangan cerminnya: $\Delta p \approx +sp$ selama
$M$ jarang (awal yang eksponensial), lalu $\Delta q \approx -sq^{2}p \to$ sebuah rangkakan ketika \omterm{def:b2:meiosis-heredity:vocabulary}{alel} yang pucat bersembunyi di dalam
\omterm{def:b2:meiosis-heredity:vocabulary}{heterozigot}. Rekursinya memberikan 27 generasi bagi kembalinya dari $p = 0.68$ menjadi $0.001$, dengan 16 generasi di antaranya untuk membawa
\omterm{def:b2:meiosis-heredity:vocabulary}{alelnya} menjadi \qty{5}{\%}.
\textbf{7.} $24$ salinan gen: $0.95^{24} = 0.29$.
\textbf{8.} $H_1 = 0.4\,(1 - 1/24) = 0.383$.
\textbf{9.} $H_{100} = 0.383\,(1 - 1/80)^{100} = 0.383\times 0.284 =
0.109$: \qty{27}{\%} dari $0.4$ daratannya.
\textbf{10.} $P(\text{fiksasi}) = 1/(2N) = 1/80 = 0.0125$; waktu
harapan sampai fiksasinya $4N = 160$ generasi.
\textbf{11.} Tanpa \omterm{prop:b2:population-genetics:inbreeding}{silang dalamnya}, $q^{2} = 0.0025$; dengan $F = 0.3$,
$q^{2} + Fpq = 0.0025 + 0.3\times 0.95\times 0.05 = 0.017$: tujuh kali
lebih banyak burung yang terkena.
\textbf{12.} Satu migran tiap generasi ($m = 1/40$) sudah cukup untuk
mencegah pulaunya memencar pada lokus netral dan untuk mengisi ulang
\omterm{def:b2:meiosis-heredity:vocabulary}{alel} yang hilang karena hanyutan; ia juga mengimpor \omterm{def:b2:meiosis-heredity:vocabulary}{alel} daratan yang
merugikan pada frekuensi daratannya, yang, di bawah \omterm{prop:b2:population-genetics:inbreeding}{silang dalam}
pulaunya, tersingkap sebagai \omterm{def:b2:meiosis-heredity:vocabulary}{homozigot} --- bebannya ditetapkan oleh
keseimbangan antara imigrasi dan penyingkapannya.
\textbf{13.} $\hat q = \sqrt{10^{-5}} = 0.0032$; kelahiran yang terkena
$\hat
q^{2} = 10^{-5}$.
\textbf{14.} $\hat q = \sqrt{10^{-5}/0.1} = 0.01$; kelahiran yang
terkena $10^{-4}$: seleksi yang sepuluh kali lebih lemah membiarkan
\omterm{def:b2:meiosis-heredity:vocabulary}{alelnya} menjadi tiga kali lebih lazim dan penyakitnya sepuluh kali.
\textbf{15.} $\hat p = \mu/(hs) = 10^{-5}/0.05 = 2\times 10^{-4}$;
pembawanya $2\hat p = 4\times 10^{-4}$.
\textbf{16.} Dari $\hat q = 0.0032$ menjadi $0.01$: \omterm{def:b2:meiosis-heredity:vocabulary}{alelnya} naik pada
sebuah laju yang ditetapkan \omterm{def:b2:genome-diversification:mutation}{mutasinya}, $\mu = 10^{-5}$ tiap generasi,
sehingga pendekatannya memerlukan orde $\Delta q/\mu \approx 700$
generasi --- \num{20000} tahun: keputusan apa pun yang dibuat sekarang
dirasakan sebuah populasi yang tidak akan menyerupai populasi kita.
\textbf{17.} Tiap lokusnya menyumbang $2\hat q = 0.0063$ salinan yang
mematikan tiap orangnya; sepanjang \num{20000} lokus, kira-kira $130$
--- tetapi itu menganggap tiap gennya dapat bermutasi menjadi sebuah
\omterm{def:b2:meiosis-heredity:vocabulary}{alel} \omterm{def:b2:meiosis-heredity:vocabulary}{resesif} yang mematikan; dengan taksiran lazim berupa beberapa ribu
gen yang hakiki, seseorang membawa beberapa sampai beberapa lusin
setara mematikan dalam keadaan \omterm{def:b2:meiosis-heredity:vocabulary}{heterozigot}.
\textbf{18.} Masing-masingnya berada pada lokus yang berbeda dan
masing-masingnya jarang, sehingga peluang bahwa sepasang acak
sama-sama membawa yang sama adalah kecil; sepupu berbagi seperdelapan
gennya menurut keturunan, sehingga bagi tiap \omterm{def:b2:meiosis-heredity:vocabulary}{alel} mematikan seorang
sepupu risiko bahwa anaknya \omterm{def:b2:meiosis-heredity:vocabulary}{homozigot} adalah $1/16$ tiap \omterm{def:b2:meiosis-heredity:vocabulary}{alel} mematikan
yang dibagi --- kelebihan kematian anak pada perkawinan sepupu terukur
sebesar beberapa persen.
\textbf{19.} $S = 1.16\times \qty{800}{L} = \qty{930}{L}$; $R = 0.4\times
930 = \qty{370}{L}$.
\textbf{20.} Lima generasi: kira-kira \qty{1860}{L}; lebih kecil dalam
praktiknya karena ragam aditifnya terkuras, karena $h^{2}$-nya ditaksir
pada kawanan aslinya, dan karena lingkungannya (pakan, kandang) yang
membuat sapi teratasnya tidak terwariskan.
\textbf{21.} Diferensial reratanya $(1.16 + 2.4)/2 = 1.78$ simpangan
baku, $S = \qty{1424}{L}$, $R = 0.4\times 1424 = \qty{570}{L}$.
\textbf{22.} Klona tidak memiliki ragam genetik, sehingga $h^{2} = 0$
dan $R =
0$ berapa pun $S$-nya: heritabilitas adalah sifat sebuah
populasi, yaitu bagian ragamnya yang bersifat genetik, bukan sifat
sifatnya.
\textbf{23.} $R = 0.7\times 0.6 = \qty{0.42}{mm}$ lebih dalam.
\textbf{24.} Tanggapan yang sama sebesar $\qty{370}{L}$ memerlukan $S = 370/0.1 =
\qty{3700}{L}$, yaitu $4.6$ simpangan baku --- dengan
menyimpan kurang dari satu sapi dari sepuluh ribu: tak praktis, dan
itulah sebabnya sifat yang heritabilitasnya rendah diperbaiki dengan
lambat atau dengan cara lain (pengujian keturunan, keluarga yang lebih
besar).
\textbf{25.} Ngengatnya: $s \approx 0.2$ sampai $0.3$, 28 generasi pada
$0.3$ dan 45 pada $0.2$, berbanding 50 yang teramati; pulaunya:
$H_{100} \approx 0.11$, yaitu \qty{27}{\%} daratannya; kawanannya:
\qty{370}{L} tiap generasi dari sapinya saja, \qty{570}{L} dengan
pejantan yang terseleksi.
\end{solution}
